\documentclass[aspectratio=169]{beamer}

\usetheme{Madrid}
\usepackage[spanish]{babel}
\usepackage[utf8]{inputenc}
\usepackage{textcomp}
\usepackage[T1]{fontenc}
\usepackage{lmodern}
\usepackage{graphicx}
\usepackage{booktabs}
\usepackage{verbatim}
\usepackage{amssymb}
\usepackage{pifont}
\usepackage{moreverb}
\usepackage[most]{tcolorbox}
\usepackage{tikz}
\usetikzlibrary{shapes, arrows, positioning}

% Configuración de TikZ
\tikzstyle{block} = [
  rectangle, draw, rounded corners,
  text centered, minimum height=1cm,
  minimum width=2.5cm, font=\small, fill=gray!5
]
\tikzstyle{arrow} = [->, thick]

% Configuración de tcolorbox
\tcbset{
  colback=gray!5, colframe=gray!50,
  boxrule=0.4pt, sharp corners,
}

\title[BookMate]{BookMate — Sistema de Recomendación de Libros Académicos}
\subtitle{Especificación de Requisitos de Software (SRS) y Casos de Uso}
\author[GRUPO 6]{Delgado R., G. \and Osorio M., A. \and Rojas A., J. \and Torres R., J. \and Valverde G., Y. \and Villanueva A., F.}
\institute[FC-UNI]{UNIVERSIDAD NACIONAL DE INGENIERÍA}
\titlegraphic{\includegraphics[height=1.4cm]{example-image}}
% NOTA: Reemplazar 'example-image' con 'Imagenes/logo_uni.png' cuando esté disponible

\date{Octubre 2025}

\begin{document}

% Portada
\begin{frame}
\titlepage
\end{frame}

% Agenda
\begin{frame}
\frametitle{Agenda}
\tableofcontents
\end{frame}

\section{Introducción al SRS}

\begin{frame}
\frametitle{¿Qué es un SRS?}

\begin{itemize}
\item \textbf{Especificación de Requisitos de Software}
\item Documento que define qué debe hacer el sistema
\item Base para el diseño y desarrollo
\item Contrato entre cliente y desarrolladores
\end{itemize}

\vspace{0.5cm}

\begin{tcolorbox}[title=Objetivo del SRS BookMate]
Definir los requisitos funcionales y no funcionales del sistema BookMate para la recomendación inteligente de libros académicos.
\end{tcolorbox}

\end{frame}

\begin{frame}
\frametitle{Descripción General del Producto}

\begin{columns}
\begin{column}{0.4\textwidth}
\textbf{Perspectiva del Producto:}
\begin{itemize}
\item Sistema web independiente
\item Arquitectura de microservicios
\item Backend + IA + Frontend
\item Base de datos
\end{itemize}

\vspace{0.3cm}

\textbf{Funciones Principales:}
\begin{itemize}
\item Gestión de catálogo (CRUD)
\item Búsqueda avanzada
\item Recomendaciones inteligentes
\item Interfaz web responsive
\end{itemize}
\end{column}

\begin{column}{0.6\textwidth}
\begin{center}
\includegraphics[width=\textwidth]{example-image-a}
% NOTA: Reemplazar con diagrama conceptual del sistema
% Mostrar: Frontend ↔ Backend ↔ IA + BD
\end{center}
\end{column}
\end{columns}

\end{frame}

\section{Requisitos Funcionales}

\begin{frame}
\frametitle{Requisitos Funcionales Principales}

\begin{table}
\centering
\small
\begin{tabular}{p{2.5cm}p{7cm}p{2cm}}
\toprule
\textbf{ID} & \textbf{Descripción} & \textbf{Prioridad} \\
\midrule
RF-01 & Gestionar Libros (CRUD) & Alta \\
RF-02 & Gestionar Autores (CRUD) & Alta \\
RF-03 & Buscar Libros & Alta \\
RF-04 & Obtener Recomendaciones & \textbf{Muy Alta} \\
RF-05 & Visualizar Detalle de Libro & Media \\
RF-06 & Filtrar por Precio & Media \\
RF-07 & Navegar Catálogo & Media \\
\bottomrule
\end{tabular}
\end{table}

\vspace{0.3cm}

\begin{tcolorbox}[colback=blue!5, colframe=blue!50]
\textbf{RF-04} es el requisito más crítico: Sistema de recomendaciones con IA que debe responder en <3 segundos.
\end{tcolorbox}

\end{frame}

\begin{frame}
\frametitle{RF-01: Gestionar Libros}

\begin{columns}
\begin{column}{0.5\textwidth}
\textbf{Descripción:}
El sistema debe permitir crear, leer, actualizar y eliminar libros.

\textbf{Entradas:}
\begin{itemize}
\item Título
\item Autor
\item Género
\item Precio
\item Fecha de edición
\item Sinopsis
\item ISBN, Editorial
\end{itemize}
\end{column}

\begin{column}{0.5\textwidth}
\textbf{Salidas:}
\begin{itemize}
\item Confirmación de operación
\item Lista de libros
\end{itemize}

\textbf{Precondiciones:}
\begin{itemize}
\item Usuario con permisos de admin
\end{itemize}

\textbf{Postcondiciones:}
\begin{itemize}
\item Libro persistido en BD
\item Embedding generado (si IA disponible)
\end{itemize}
\end{column}
\end{columns}

\end{frame}

\begin{frame}
\frametitle{RF-04: Obtener Recomendaciones (Crítico)}

\begin{columns}
\begin{column}{0.5\textwidth}
\textbf{Descripción:}
El sistema debe recomendar libros similares basándose en embeddings semánticos o heurísticas.

\textbf{Entradas:}
\begin{itemize}
\item ID de libro de referencia
\end{itemize}

\textbf{Salidas:}
\begin{itemize}
\item Lista de 6 libros recomendados
\item Score de similitud
\end{itemize}
\end{column}

\begin{column}{0.5\textwidth}
\textbf{Procesamiento:}
\begin{enumerate}
\item Obtener datos del libro
\item Enviar a servicio IA
\item IA calcula similitud
\item Retornar top 6
\end{enumerate}

\textbf{Fallback:}
\begin{itemize}
\item Si IA no disponible
\item Aplicar heurísticas
\item Basadas en género, autor, tags
\end{itemize}
\end{column}
\end{columns}

\vspace{0.3cm}

\begin{tcolorbox}[colback=green!5, colframe=green!50]
\textbf{Objetivo:} >70\% relevancia, <3s respuesta
\end{tcolorbox}

\end{frame}

\section{Requisitos No Funcionales}

\begin{frame}
\frametitle{Requisitos No Funcionales}

\begin{table}
\centering
\small
\begin{tabular}{p{3cm}p{6cm}p{2cm}}
\toprule
\textbf{Categoría} & \textbf{Requisito} & \textbf{Objetivo} \\
\midrule
Rendimiento & Búsquedas & <1s (p95) \\
Rendimiento & Recomendaciones IA & <3s (p95) \\
Rendimiento & API REST & <500ms (p95) \\
Usabilidad & Interfaz intuitiva & <5 clicks \\
Usabilidad & Responsive design & Móvil/Tablet/Desktop \\
Fiabilidad & Disponibilidad & 95\% \\
Seguridad & Validación entrada & 100\% \\
Mantenibilidad & Cobertura tests & >80\% \\
Escalabilidad & Usuarios concurrentes & 100 \\
\bottomrule
\end{tabular}
\end{table}

\end{frame}

\begin{frame}
\frametitle{Atributos de Calidad}

\begin{table}
\centering
\begin{tabular}{p{4cm}p{3cm}p{2cm}}
\toprule
\textbf{Atributo} & \textbf{Métrica} & \textbf{Objetivo} \\
\midrule
Precisión recomendaciones & \% relevancia & >70\% \\
Tiempo respuesta búsqueda & Latencia p95 & <1s \\
Tiempo respuesta IA & Latencia p95 & <3s \\
Cobertura tests & \% código & >80\% \\
Uptime & \% disponibilidad & 95\% \\
\bottomrule
\end{tabular}
\end{table}

\vspace{0.3cm}

\begin{tcolorbox}[colback=yellow!5, colframe=yellow!50]
\textbf{Objetivo Principal:} Sistema de recomendaciones con >70\% de precisión y respuesta en <3 segundos.
\end{tcolorbox}

\end{frame}

\section{Casos de Uso Principales}

\begin{frame}
\frametitle{Actores del Sistema}

\begin{columns}
\begin{column}{0.5\textwidth}
\textbf{Actor Principal: Usuario}
\begin{itemize}
\item Estudiante
\item Profesor
\item Investigador
\end{itemize}

\textbf{Casos de Uso:} UC-02, UC-05, UC-06, UC-07, UC-08, UC-09
\end{column}

\begin{column}{0.5\textwidth}
\textbf{Actor Secundario: Administrador}
\begin{itemize}
\item Bibliotecario
\item Gestor de catálogo
\end{itemize}

\textbf{Casos de Uso:} UC-01, UC-03, UC-04, UC-10 a UC-15

\vspace{0.3cm}

\textbf{Actor Externo: Sistema IA}
\begin{itemize}
\item Servicio Python
\item Genera recomendaciones
\end{itemize}
\end{column}
\end{columns}

\vspace{0.5cm}

\begin{center}
\begin{tikzpicture}[node distance=2cm]
\node[block] (usuario) {Usuario};
\node[block, right=of usuario] (sistema) {Sistema BookMate};
\node[block, right=of sistema] (ia) {Servicio IA};
\draw[arrow] (usuario) -- (sistema);
\draw[arrow] (sistema) -- (ia);
\end{tikzpicture}
\end{center}

\end{frame}

\begin{frame}
\frametitle{UC-01: Crear Libro}

\textbf{Actor:} Administrador

\textbf{Precondiciones:} Autor existe o se creará

\textbf{Flujo Principal:}
\begin{enumerate}
\item Usuario selecciona ``Agregar Libro''
\item Sistema muestra formulario
\item Usuario ingresa: título, autor, género, precio, fecha, sinopsis, etc.
\item Sistema valida datos
\item Sistema guarda libro en BD
\item Sistema notifica a servicio IA para generar embedding
\item Sistema muestra confirmación
\end{enumerate}

\textbf{Flujo Alternativo 4a:} Datos inválidos
\begin{itemize}
\item Sistema muestra errores específicos
\item Retorna al paso 3
\end{itemize}

\end{frame}

\begin{frame}
\frametitle{UC-06: Obtener Recomendaciones (Crítico)}

\textbf{Actor Principal:} Usuario  
\textbf{Actores Secundarios:} Servicio IA

\textbf{Precondiciones:} Libro de referencia existe

\textbf{Flujo Principal:}
\begin{enumerate}
\item Usuario selecciona ``Ver similares'' en libro X
\item Sistema verifica disponibilidad de Servicio IA
\item Sistema envía solicitud a Servicio IA con ID de libro X
\item Servicio IA genera embedding de sinopsis
\item Servicio IA calcula similitud de coseno con todos los libros
\item Servicio IA retorna lista de IDs de libros recomendados
\item Sistema consulta detalles en BD
\item Sistema muestra recomendaciones con scores
\end{enumerate}

\textbf{Flujo Alternativo 2a:} Servicio IA no disponible
\begin{itemize}
\item Sistema aplica heurísticas (mismo género/autor/precio/tags)
\item Sistema muestra recomendaciones heurísticas
\end{itemize}

\end{frame}

\begin{frame}
\frametitle{UC-05: Buscar Libros}

\textbf{Actor:} Usuario

\textbf{Precondiciones:} Ninguna

\textbf{Flujo Principal:}
\begin{enumerate}
\item Usuario ingresa query de búsqueda
\item Sistema busca coincidencias en título/autor/género
\item Sistema ordena por relevancia
\item Sistema muestra resultados
\end{enumerate}

\textbf{Flujo Alternativo 3a:} Sin resultados
\begin{itemize}
\item Sistema muestra mensaje ``Sin resultados''
\item Sistema sugiere libros populares
\end{itemize}

\vspace{0.3cm}

\begin{tcolorbox}[colback=blue!5, colframe=blue!50]
\textbf{Objetivo:} Búsquedas deben responder en <1 segundo (RNF-01.1)
\end{tcolorbox}

\end{frame}

\section{Diagrama de Casos de Uso}

\begin{frame}
\frametitle{Diagrama de Casos de Uso}

\begin{columns}
\begin{column}{0.4\textwidth}
\begin{itemize}
\item \textbf{Actores:}
  \begin{itemize}
  \item Usuario
  \item Administrador
  \item Servicio IA
  \end{itemize}
\item \textbf{Casos de Uso Principales:}
  \begin{itemize}
  \item UC-01: Crear Libro
  \item UC-02: Listar Libros
  \item UC-05: Buscar Libros
  \item UC-06: Obtener Recomendaciones
  \item UC-10-13: Gestionar Autores
  \end{itemize}
\end{itemize}
\end{column}

\begin{column}{0.6\textwidth}
\begin{center}
\includegraphics[width=\textwidth]{example-image-b}
% NOTA: Reemplazar con diagrama casos_uso.puml generado
% Mostrar el diagrama principal con todos los actores y UC
\end{center}
\end{column}
\end{columns}

\end{frame}

\begin{frame}
\frametitle{Diagrama Detallado: UC-06 Recomendaciones}

\begin{center}
\includegraphics[width=0.9\textwidth]{example-image-c}
% NOTA: Reemplazar con el diagrama detallado de UC-06
% Mostrar flujo principal (IA) y flujo alternativo (heurísticas)
\end{center}

\textbf{Nota:} Este es el caso de uso más crítico del sistema. Define el valor principal del producto.

\end{frame}

\section{Matriz de Trazabilidad}

\begin{frame}
\frametitle{Matriz de Trazabilidad: Requisitos - Casos de Uso}

\begin{table}
\centering
\small
\begin{tabular}{p{2cm}p{4cm}p{4cm}p{2cm}}
\toprule
\textbf{Requisito} & \textbf{Descripción} & \textbf{Casos de Uso} & \textbf{Prioridad} \\
\midrule
RF-01 & Gestionar Libros & UC-01, UC-02, UC-03, UC-04 & Alta \\
RF-02 & Gestionar Autores & UC-10, UC-11, UC-12, UC-13 & Alta \\
RF-03 & Buscar Libros & UC-05 & Alta \\
RF-04 & Obtener Recomendaciones & UC-06 & \textbf{Muy Alta} \\
RF-05 & Visualizar Detalle & UC-07 & Media \\
RF-06 & Filtrar por Precio & UC-08 & Media \\
RF-07 & Navegar Catálogo & UC-09 & Media \\
\bottomrule
\end{tabular}
\end{table}

\end{frame}

\begin{frame}
\frametitle{Matriz de Trazabilidad: Casos de Uso - Componentes}

\begin{table}
\centering
\tiny
\begin{tabular}{p{2.5cm}p{10cm}}
\toprule
\textbf{Caso de Uso} & \textbf{Componentes del Sistema (Arquitectura Futura)} \\
\midrule
UC-01 & LibrosController, LibrosService, LibrosRepository, Libros, Autor \\
UC-02 & LibrosController, LibrosService, LibrosRepository, Libros \\
UC-03 & LibrosController, LibrosService, LibrosRepository, Libros \\
UC-04 & LibrosController, LibrosService, LibrosRepository \\
UC-05 & LibrosController, WebController, LibrosService, LibrosRepository \\
UC-06 & RecommendationController, RecommendationService, AIService, HeuristicRecommender \\
UC-07 & WebController, LibrosService, Libros \\
UC-08 & LibrosController, LibrosService, LibrosRepository \\
UC-10-13 & AutorController, AutorService, AutorRepository, Autor \\
\bottomrule
\end{tabular}
\end{table}

\vspace{0.3cm}

\textbf{Nota:} Esta es la arquitectura planificada. En la fase de análisis (actual) no se mencionan tecnologías específicas.

\end{frame}

\begin{frame}
\frametitle{Matriz de Trazabilidad: Casos de Uso - Pruebas}

\begin{table}
\centering
\tiny
\begin{tabular}{p{1.5cm}p{5.5cm}p{5.5cm}p{1cm}}
\toprule
\textbf{Caso de Uso} & \textbf{Pruebas Unitarias (Futuras)} & \textbf{Pruebas API (Futuras)} & \textbf{Estado} \\
\midrule
UC-01 & LibrosServiceTest.testSave & POST /api/libros & Planificado \\
UC-02 & LibrosServiceTest.testFindAll & GET /api/libros & Planificado \\
UC-03 & LibrosServiceTest.testUpdate & PUT /api/libros/\{id\} & Planificado \\
UC-04 & LibrosServiceTest.testDelete & DELETE /api/libros/\{id\} & Planificado \\
UC-05 & LibrosServiceTest.testSearchBooks & GET /api/libros/search & Planificado \\
UC-06 & RecommendationServiceTest & GET /api/recommendations/\{id\} & Planificado \\
UC-08 & LibrosServiceTest.testFindByPrecio & GET /api/libros/precio & Planificado \\
\bottomrule
\end{tabular}
\end{table}

\end{frame}

\section{Prototipo de Casos de Uso}

\begin{frame}
\frametitle{Prototipo Actual}

\textbf{Estado:} Funcional

\textbf{Ubicación:} \texttt{product/basic-springboot/}

\begin{columns}
\begin{column}{0.5\textwidth}
\textbf{Funcionalidades Implementadas:}
\begin{itemize}
\item[\checkmark] Interfaz web responsive
\item[\checkmark] Catálogo navegable
\item[\checkmark] Búsqueda básica
\item[\checkmark] Vista de detalles
\item[\checkmark] Recomendaciones heurísticas
\item[\checkmark] Panel de admin (CRUD)
\end{itemize}
\end{column}

\begin{column}{0.5\textwidth}
\textbf{Limitaciones:}
\begin{itemize}
\item[\ding{55}] Sin backend real
\item[\ding{55}] Sin servicio IA
\item[\ding{55}] Datos en LocalStorage
\item[\ding{55}] Sin gestión de autores
\item[\ding{55}] Sin importar/exportar CSV
\end{itemize}
\end{column}
\end{columns}

\vspace{0.3cm}

\begin{tcolorbox}[colback=green!5, colframe=green!50]
\textbf{Propósito:} Validar requisitos funcionales y obtener feedback temprano
\end{tcolorbox}

\end{frame}

\begin{frame}
\frametitle{Capturas del Prototipo}

\begin{center}
\includegraphics[width=0.45\textwidth]{example-image}
\includegraphics[width=0.45\textwidth]{example-image-a}
% NOTA: Reemplazar con capturas reales del prototipo
% Izquierda: Catálogo de libros
% Derecha: Panel de administración
\end{center}

\begin{center}
\includegraphics[width=0.45\textwidth]{example-image-b}
\includegraphics[width=0.45\textwidth]{example-image-c}
% NOTA: Reemplazar con capturas reales
% Izquierda: Detalle de libro
% Derecha: Recomendaciones
\end{center}

\end{frame}

\begin{frame}
\frametitle{Flujo de Usuario en el Prototipo}

\textbf{Escenario:} Usuario descubre libros relacionados

\begin{enumerate}
\item Usuario accede al catálogo (\texttt{/catalog.html})
\item Explora libros o usa búsqueda
\item Hace click en un libro de interés
\item Visualiza detalles completos
\item Hace click en ``Ver libros similares''
\item Sistema muestra 6 recomendaciones heurísticas:
  \begin{itemize}
  \item Mismo género: +3 puntos
  \item Autor común: +5 puntos
  \item Tags compartidos: +2 puntos cada uno
  \item Precio similar: +1 punto
  \end{itemize}
\item Usuario explora las recomendaciones
\end{enumerate}

\vspace{0.3cm}

\textbf{Resultado:} Usuario descubre libros relevantes en <2 minutos

\end{frame}

\section{Criterios de Aceptación}

\begin{frame}
\frametitle{Criterios de Aceptación}

\begin{columns}
\begin{column}{0.5\textwidth}
\textbf{Requisitos Funcionales:}
\begin{itemize}
\item [\checkmark] CRUD completo de libros
\item [\checkmark] CRUD completo de autores
\item [\checkmark] Búsqueda por múltiples criterios
\item [\checkmark] Sistema de recomendaciones IA
\item [\checkmark] Sistema de recomendaciones heurísticas
\item [\checkmark] Filtros por precio y género
\end{itemize}
\end{column}

\begin{column}{0.5\textwidth}
\textbf{Requisitos No Funcionales:}
\begin{itemize}
\item [\checkmark] Búsquedas <1 segundo
\item [\checkmark] Recomendaciones IA <3 segundos
\item [\checkmark] Interfaz responsive
\item [\checkmark] Cobertura tests >80\%
\item [\checkmark] Disponibilidad 95\%
\item [\checkmark] Relevancia recomendaciones >70\%
\end{itemize}
\end{column}
\end{columns}

\vspace{0.5cm}

\begin{tcolorbox}[colback=green!5, colframe=green!50]
\textbf{Objetivo Principal:} Sistema de recomendaciones con >70\% de precisión y respuesta en <3 segundos.
\end{tcolorbox}

\end{frame}

\begin{frame}
\frametitle{Restricciones del Proyecto}

\begin{itemize}
\item \textbf{Tiempo:} 10 semanas de desarrollo (S7-S16)
\item \textbf{Equipo:} 6 integrantes con disponibilidad limitada
\item \textbf{Recursos:} Sin presupuesto para servicios cloud
\item \textbf{Fases:}
  \begin{itemize}
  \item \textbf{Análisis (S7-S11):} Sin mencionar tecnologías específicas
  \item \textbf{Diseño (S13-S16):} Con implementación técnica detallada
  \end{itemize}
\end{itemize}

\vspace{0.3cm}

\begin{tcolorbox}[colback=red!5, colframe=red!50]
\textbf{Interfaz de Software (Planificado para fase de diseño):}
\begin{itemize}
\item Backend: Framework web moderno
\item BD: Sistema relacional
\item IA: Microservicio Python
\item Comunicación: HTTP/REST
\end{itemize}
\end{tcolorbox}

\textbf{Nota:} Las tecnologías específicas se definirán en la Semana 13 (fase de diseño).

\end{frame}

\section{Resumen}

\begin{frame}
\frametitle{Resumen del SRS}

\begin{columns}
\begin{column}{0.5\textwidth}
\textbf{Requisitos Funcionales:}
\begin{itemize}
\item 15 requisitos definidos
\item RF-04 (Recomendaciones) es crítico
\item 15 casos de uso especificados
\item Matriz de trazabilidad completa
\end{itemize}
\end{column}

\begin{column}{0.5\textwidth}
\textbf{Requisitos No Funcionales:}
\begin{itemize}
\item Rendimiento: <1s búsqueda, <3s IA
\item Usabilidad: <5 clicks, responsive
\item Fiabilidad: 95\% disponibilidad
\item Mantenibilidad: >80\% cobertura
\item Precisión: >70\% relevancia
\end{itemize}
\end{column}
\end{columns}

\vspace{0.5cm}

\begin{tcolorbox}[colback=blue!5, colframe=blue!50, title=Objetivo Principal]
Sistema de recomendación de libros académicos con IA que responda en <3 segundos y tenga >70\% de precisión.
\end{tcolorbox}

\end{frame}

\begin{frame}
\frametitle{Cobertura de Requisitos}

\textbf{Trazabilidad:}

\begin{itemize}
\item ✅ \textbf{100\%} de requisitos funcionales cubiertos por casos de uso
\item ✅ \textbf{100\%} de requisitos no funcionales mapeados
\item ✅ Todos los casos de uso tienen componentes identificados (arquitectura futura)
\item ✅ Todos los casos de uso tienen pruebas planificadas
\item ✅ Prototipo funcional valida requisitos principales
\end{itemize}

\vspace{0.5cm}

\textbf{Casos de Uso Críticos:}
\begin{enumerate}
\item \textbf{UC-06: Obtener Recomendaciones} (Muy Alta prioridad)
\item UC-01: Crear Libro (Alta prioridad)
\item UC-05: Buscar Libros (Alta prioridad)
\end{enumerate}

\end{frame}

\begin{frame}
\frametitle{Próximos Pasos}

\textbf{Semana 11 - Análisis Conceptual:}

\begin{itemize}
\item Especificaciones detalladas de casos de uso (UC-01, UC-05, UC-06)
\item Modelo de análisis (sin tecnologías específicas)
\item Diagramas de robustez (Boundary-Control-Entity)
\item Diagramas de secuencia conceptuales
\item Arquitectura C4 Niveles 1 y 2 (Contexto y Contenedores)
\item Tarjetas CRC (Class-Responsibility-Collaboration)
\item Diagramas de estado (si aplica)
\end{itemize}

\vspace{0.3cm}

\begin{tcolorbox}[colback=yellow!5, colframe=yellow!50]
\textbf{Importante:} Continuar en fase de análisis sin mencionar tecnologías hasta Semana 13
\end{tcolorbox}

\end{frame}

\begin{frame}
\frametitle{¡Preguntas?}

\begin{center}
\Large ¡Gracias por su atención! \\

\vspace{1cm}

\textbf{¿Preguntas?}

\vspace{1cm}

\normalsize
\textbf{Grupo 6}\\
Universidad Nacional de Ingeniería\\
CC341 - Ingeniería de Software\\
Ciclo 2025-2

\end{center}

\end{frame}

\end{document}

